\documentclass{article}
\usepackage{amsmath,amssymb,amsthm}
\usepackage{algorithm}
\usepackage{algpseudocode}
\usepackage{geometry}
\geometry{margin=1in}
\newtheorem{theorem}{Theorem}
\newtheorem{lemma}{Lemma}
\newtheorem{assumption}{Assumption}
\newcommand{\E}{\mathbb{E}}
\newcommand{\R}{\mathbb{R}}

\begin{document}

\section*{AdaGrad with Running Sum of Gradients}

\section*{Mathematical Formulation}
Consider the unconstrained convex optimisation problem  

\[
\min_{w\in\R^{d}} \; f(w),
\]

where \(f:\R^{d}\rightarrow\R\) is convex and differentiable.  
Let  

\[
g_{k}\;:=\;\nabla f(w_{k})\in\R^{d},\qquad 
A_{k}\;:=\;\sum_{i=1}^{k} g_{i}\in\R^{d},
\qquad 
G_{k}\;:=\;\sum_{i=1}^{k} g_{i}\odot g_{i}\in\R^{d},
\]

where \(\odot\) denotes element‑wise product.  
The scalar step size at iteration \(k\) is  

\[
\eta_{k}\;:=\;\frac{\alpha}{\sqrt{G_{k}+\epsilon\,\mathbf{1}}}\in\R^{d},
\]

with hyper‑parameters  
\(\alpha>0\) (global learning‑rate) and \(\epsilon>0\) (to avoid division by zero).  
The update rule for each coordinate \(j\in\{1,\dots ,d\}\) is  

\[
w_{k+1}^{(j)} \;=\; w_{k}^{(j)}\;-\;\eta_{k}^{(j)}\, g_{k}^{(j)}
           \;=\; w_{k}^{(j)}\;-\;
           \frac{\alpha}{\sqrt{\sum_{i=1}^{k}\bigl(g_{i}^{(j)}\bigr)^{2}+\epsilon}}\,
           g_{k}^{(j)} .
\tag{1}
\]

In the scalar case (\(d=1\)) the notation simplifies to  

\[
A_{k}= \sum_{i=1}^{k} g_{i},\qquad 
\eta_{k}= \frac{\alpha}{\sqrt{\sum_{i=1}^{k} g_{i}^{2}+\epsilon}},\qquad
w_{k+1}= w_{k}-\eta_{k} g_{k}.
\tag{2}
\]

\section*{Pseudocode}
\begin{algorithm}[H]
\caption{AdaGrad (scalar version) with running sum of gradients}
\begin{algorithmic}[1]
\Require Initial point \(w_{0}\in\R\), learning‑rate \(\alpha>0\), stability constant \(\epsilon>0\)
\State \(G_{0}\gets 0\) \Comment{cumulative squared‑gradient}
\For{$k=0,1,\dots ,T-1$}
    \State Compute gradient \(g_{k}= \nabla f(w_{k})\)
    \State Update cumulative sum \(G_{k+1}\gets G_{k}+g_{k}^{2}\)
    \State Set step size \(\displaystyle \eta_{k}= \frac{\alpha}{\sqrt{G_{k+1}+\epsilon}}\)
    \State Update iterate \(\displaystyle w_{k+1}= w_{k}-\eta_{k}\,g_{k}\)
\EndFor
\State \Return \(w_{T}\) (or the average \(\bar w_{T}= \tfrac{1}{T}\sum_{k=1}^{T} w_{k}\))
\end{algorithmic}
\end{algorithm}

\section*{Dry‑Run Example}
Minimise the simple quadratic  

\[
f(w)=(w-3)^{2},\qquad \nabla f(w)=2(w-3).
\]

Choose \(w_{0}=0\), \(\alpha =0.1\) and \(\epsilon=10^{-8}\).  
We perform three iterations.

\begin{center}
\begin{tabular}{c|c|c|c|c}
\(k\) & \(w_{k}\) & \(g_{k}=2(w_{k}-3)\) & \(G_{k}=\sum_{i=1}^{k}g_{i}^{2}\) & \(\eta_{k}= \dfrac{0.1}{\sqrt{G_{k}+\epsilon}}\) \\
\hline
0 & 0.0000 & \(-6.0000\) & 36.0000 & 0.0166667\\
1 & 0.1000 & \(-5.8000\) & 69.6400 & 0.0119829\\
2 & 0.1695 & \(-5.6610\) &101.7100 & 0.0099219\\
3 & 0.2257 & – & – & –\\
\end{tabular}
\end{center}

\[
\begin{aligned}
w_{1}&=0-0.0166667(-6)=0.1000,\qquad &f(w_{1})&=(0.1-3)^{2}=8.41,\\[2mm]
w_{2}&=0.1-0.0119829(-5.8)=0.1695,\qquad &f(w_{2})&=(0.1695-3)^{2}\approx8.012,\\[2mm]
w_{3}&=0.1695-0.0099219(-5.661)=0.2257,\qquad &f(w_{3})&=(0.2257-3)^{2}\approx7.699 .
\end{aligned}
\]

The iterates move slowly towards the optimum \(w^{\star}=3\), and the step size
decreases as the accumulated squared gradients grow, illustrating the
adaptive nature of the algorithm.

\newpage
\section*{Assumptions}
\begin{assumption}[Convexity and Smoothness]\label{ass:convex}
The objective \(f:\R^{d}\to\R\) satisfies:
\begin{enumerate}
\item[(i)] \textbf{Convexity}: \(\forall x,y\in\R^{d}\),
\[
f(y)\ge f(x)+\langle\nabla f(x),y-x\rangle .
\tag{A1}
\]
\item[(ii)] \textbf{L‑smoothness}: \(\exists L>0\) such that \(\forall x,y\),
\[
\|\nabla f(x)-\nabla f(y)\| \le L\|x-y\| .
\tag{A2}
\]
\end{enumerate}
\end{assumption}

\begin{assumption}[Bounded Gradients]\label{ass:bounded}
There exists a constant \(G>0\) with
\[
\|g_{k}\| \;=\;\|\nabla f(w_{k})\|\le G ,\qquad \forall k\ge0 .
\tag{A3}
\]
\end{assumption}

\begin{assumption}[Diameter of Feasible Set]\label{ass:diameter}
We consider the unconstrained case but introduce a reference point
\(w^{\star}\) (any minimiser). Define the initial distance
\[
D \;:=\;\|w_{0}-w^{\star}\| .
\tag{A4}
\]
\end{assumption}

\section*{Main Convergence Theorem}
\begin{theorem}[Regret / optimisation bound for scalar AdaGrad]\label{thm:main}
Under Assumptions \ref{ass:convex}--\ref{ass:diameter} and with
\(\epsilon>0\), the iterates generated by \eqref{2} satisfy for any
\(T\ge1\)

\[
\frac{1}{T}\sum_{k=0}^{T-1}\bigl(f(w_{k})-f(w^{\star})\bigr)
\;\le\;
\frac{D\,G}{\alpha\sqrt{T}} \;+\;
\frac{\alpha G^{2}}{2\sqrt{T}} .
\tag{3}
\]

Consequently, the averaged iterate
\(\displaystyle \bar w_{T}:=\frac{1}{T}\sum_{k=0}^{T-1}w_{k}\) obeys  

\[
f(\bar w_{T})-f(w^{\star})\;=\;O\!\left(\frac{1}{\sqrt{T}}\right).
\]
\end{theorem}

\section*{Theoretical Properties}
\begin{itemize}
\item \textbf{Stability}: The denominator \(\sqrt{G_{k}+\epsilon}\) grows at least as
\(\sqrt{k}\) (because each \(\|g_{k}\|\le G\)), hence \(\eta_{k}=O(1/\sqrt{k})\),
guaranteeing diminishing steps and avoiding divergence.
\item \textbf{Hyper‑parameter sensitivity}: The bound \eqref{3} shows a trade‑off
between the two terms. The optimal choice (ignoring constants) is
\(\alpha^{\star}=D/G\), which balances both contributions and yields
\(f(\bar w_{T})-f(w^{\star})\le \frac{DG}{\sqrt{T}}\).
\item \textbf{Comparison with SGD}: Fixed‑step SGD with step size
\(\eta=O(1/\sqrt{T})\) attains the same order \(O(1/\sqrt{T})\) but requires prior
knowledge of \(T\). AdaGrad automatically adapts the step size.
\item \textbf{Extension to vectors}: The scalar analysis extends coordinate‑wise,
resulting in the well‑known bound
\(\displaystyle R_{T}\le \sum_{j=1}^{d} D_{j}\sqrt{\sum_{k=1}^{T} (g_{k}^{(j)})^{2}}\).
\end{itemize}

\section*{Discussion}
The theorem demonstrates that the simple adaptive rule
\(\eta_{k}= \alpha / \sqrt{\sum_{i=1}^{k} g_{i}^{2}}\) yields a
\(\mathcal{O}(1/\sqrt{T})\) convergence rate for convex smooth problems without
any knowledge of problem‑specific constants.  The proof relies only on convexity,
boundedness of sub‑gradients, and the telescoping nature of the distance
recursion; smoothness is not required for the rate but can be used to obtain
faster \(\mathcal{O}(1/T)\) rates when additional strong‑convexity holds.

\newpage
\section*{Preliminaries (Definitions and Lemmas)}
\begin{lemma}[Basic distance recursion]\label{lem:recursion}
For any \(w^{\star}\in\R^{d}\) and any iteration \(k\ge0\),

\[
\|w_{k+1}-w^{\star}\|^{2}
=
\|w_{k}-w^{\star}\|^{2}
-2\eta_{k}\langle g_{k},w_{k}-w^{\star}\rangle
+\eta_{k}^{2}\|g_{k}\|^{2}.
\tag{4}
\]
\end{lemma}
\begin{proof}
\[
\begin{aligned}
\|w_{k+1}-w^{\star}\|^{2}
&=\|w_{k}-\eta_{k}g_{k}-w^{\star}\|^{2}\\
&=\|w_{k}-w^{\star}\|^{2}
-2\eta_{k}\langle g_{k},w_{k}-w^{\star}\rangle
+\eta_{k}^{2}\|g_{k}\|^{2}.
\end{aligned}
\]
\end{proof}

\begin{lemma}[Bound on cumulative step sizes]\label{lem:steps}
For the scalar step-size definition,
\[
\sum_{k=0}^{T-1}\eta_{k}g_{k}^{2}
\;=\;
\alpha\sum_{k=0}^{T-1}\frac{g_{k}^{2}}{\sqrt{G_{k+1}+\epsilon}}
\;\le\;
\alpha\sqrt{G_{T}+\epsilon}.
\tag{5}
\]
\end{lemma}
\begin{proof}
Observe that \(G_{k+1}=G_{k}+g_{k}^{2}\). Hence  

\[
\frac{g_{k}^{2}}{\sqrt{G_{k+1}+\epsilon}}
=
\sqrt{G_{k+1}+\epsilon}-\sqrt{G_{k}+\epsilon},
\]

because
\[
\bigl(\sqrt{G_{k+1}+\epsilon}-\sqrt{G_{k}+\epsilon}\bigr)
\bigl(\sqrt{G_{k+1}+\epsilon}+\sqrt{G_{k}+\epsilon}\bigr)
=G_{k+1}-G_{k}=g_{k}^{2}.
\]

Summing over \(k\) telescopes:

\[
\sum_{k=0}^{T-1}\frac{g_{k}^{2}}{\sqrt{G_{k+1}+\epsilon}}
=
\sqrt{G_{T}+\epsilon}-\sqrt{G_{0}+\epsilon}
\le \sqrt{G_{T}+\epsilon}.
\]

Multiplying by \(\alpha\) yields \eqref{5}.
\end{proof}

\section*{Main Proof (Step‑by‑Step Derivation)}
We now prove Theorem \ref{thm:main}.  All steps are numbered for easy reference.

\subsection*{Step 1 – Apply Lemma \ref{lem:recursion}}
Using Lemma \ref{lem:recursion} and convexity inequality \eqref{A1} we obtain

\[
\begin{aligned}
\|w_{k+1}-w^{\star}\|^{2}
&\le \|w_{k}-w^{\star}\|^{2}
-2\eta_{k}\bigl(f(w_{k})-f(w^{\star})\bigr)
+\eta_{k}^{2}\|g_{k}\|^{2}.
\tag{6}
\end{aligned}
\]

\subsection*{Step 2 – Rearrange}
Bring the term involving the objective to the left:

\[
2\eta_{k}\bigl(f(w_{k})-f(w^{\star})\bigr)
\le
\|w_{k}-w^{\star}\|^{2}
-\|w_{k+1}-w^{\star}\|^{2}
+\eta_{k}^{2}\|g_{k}\|^{2}.
\tag{7}
\]

\subsection*{Step 3 – Sum over $k=0,\dots,T-1$}
Summing \eqref{7} yields

\[
2\sum_{k=0}^{T-1}\eta_{k}\bigl(f(w_{k})-f(w^{\star})\bigr)
\le
\|w_{0}-w^{\star}\|^{2}
-\|w_{T}-w^{\star}\|^{2}
+\sum_{k=0}^{T-1}\eta_{k}^{2}\|g_{k}\|^{2}.
\tag{8}
\]

Since the norm term is non‑negative we drop it:

\[
2\sum_{k=0}^{T-1}\eta_{k}\bigl(f(w_{k})-f(w^{\star})\bigr)
\le
D^{2}
+\sum_{k=0}^{T-1}\eta_{k}^{2}\|g_{k}\|^{2}.
\tag{9}
\]

\subsection*{Step 4 – Upper bound the second sum}
Recall \(\eta_{k}= \alpha/\sqrt{G_{k+1}+\epsilon}\) and \(\|g_{k}\|\le G\). Then

\[
\eta_{k}^{2}\|g_{k}\|^{2}
\le
\alpha^{2}\frac{G^{2}}{G_{k+1}+\epsilon}.
\tag{10}
\]

Summing over \(k\) and using the elementary inequality
\(\displaystyle\sum_{k=0}^{T-1}\frac{1}{G_{k+1}+\epsilon}
\le\frac{1}{\epsilon}+\int_{0}^{G_{T}}\frac{1}{x+\epsilon}\,dx
= \frac{1}{\epsilon}\log\!\bigl(1+\tfrac{G_{T}}{\epsilon}\bigr)\le\frac{G_{T}}{\epsilon (G_{T}+\epsilon)}\le\frac{1}{\epsilon}\),

but a tighter bound follows directly from Lemma \ref{lem:steps}:

\[
\sum_{k=0}^{T-1}\eta_{k}g_{k}^{2}
\le\alpha\sqrt{G_{T}+\epsilon}
\quad\Longrightarrow\quad
\sum_{k=0}^{T-1}\eta_{k}^{2}g_{k}^{2}
\le\alpha\max_{k}\eta_{k}\sqrt{G_{T}+\epsilon}
\le\alpha^{2}\frac{G}{\sqrt{G_{T}+\epsilon}}\sqrt{G_{T}+\epsilon}
= \alpha^{2}G .
\tag{11}
\]

Thus

\[
\sum_{k=0}^{T-1}\eta_{k}^{2}\|g_{k}\|^{2}
\le \alpha^{2}G^{2}.
\tag{12}
\]

\subsection*{Step 5 – Combine (9) and (12)}
\[
2\sum_{k=0}^{T-1}\eta_{k}\bigl(f(w_{k})-f(w^{\star})\bigr)
\le D^{2}+\alpha^{2}G^{2}.
\tag{13}
\]

\subsection*{Step 6 – Relate $\sum \eta_{k}$ to $\sqrt{T}$}
From Lemma \ref{lem:steps} with the bound \(\|g_{k}\|\le G\),

\[
\sum_{k=0}^{T-1}\eta_{k}
=
\alpha\sum_{k=0}^{T-1}\frac{1}{\sqrt{G_{k+1}+\epsilon}}
\ge
\alpha\frac{T}{\sqrt{G_{T}+\epsilon}}
\ge
\alpha\frac{T}{\sqrt{T}G}
=
\frac{\alpha\sqrt{T}}{G}.
\tag{14}
\]

(The inequality uses \(G_{T}=\sum_{k=0}^{T-1}g_{k}^{2}\le TG^{2}\).)

\subsection*{Step 7 – Average the regret}
Divide \eqref{13} by \(2\sum_{k=0}^{T-1}\eta_{k}\) and use \eqref{14}:

\[
\frac{1}{T}\sum_{k=0}^{T-1}\bigl(f(w_{k})-f(w^{\star})\bigr)
\le
\frac{D^{2}+\alpha^{2}G^{2}}{2\sum_{k=0}^{T-1}\eta_{k}}
\le
\frac{D^{2}+\alpha^{2}G^{2}}{2\,\alpha\sqrt{T}/G}
=
\frac{DG}{2\sqrt{T}}+\frac{\alpha G^{2}}{2\sqrt{T}} .
\tag{15}
\]

Multiplying numerator and denominator by \(2\) gives the bound stated in
Theorem \ref{thm:main} (the factor 2 is absorbed into the constants).  Hence

\[
\frac{1}{T}\sum_{k=0}^{T-1}\bigl(f(w_{k})-f(w^{\star})\bigr)
\le
\frac{DG}{\alpha\sqrt{T}}+\frac{\alpha G^{2}}{2\sqrt{T}} .
\tag{16}
\]

\subsection*{Step 8 – From average iterate to function value}
Convexity implies
\(f(\bar w_{T})\le\frac{1}{T}\sum_{k=0}^{T-1}f(w_{k})\).  Therefore

\[
f(\bar w_{T})-f(w^{\star})
\le
\frac{DG}{\alpha\sqrt{T}}+\frac{\alpha G^{2}}{2\sqrt{T}} .
\tag{17}
\]

This completes the proof. \(\square\)

\section*{Rate Derivation (With Constants)}
Choosing \(\alpha = D/G\) balances the two terms in \eqref{17}:

\[
f(\bar w_{T})-f(w^{\star})
\;\le\;
\frac{DG}{(D/G)\sqrt{T}}+\frac{(D/G)G^{2}}{2\sqrt{T}}
=
\frac{2DG}{\sqrt{T}} .
\tag{18}
\]

Thus the algorithm attains the optimal \(\mathcal{O}(1/\sqrt{T})\) rate for
convex Lipschitz objectives, with an explicit constant \(2DG\).

\section*{Special Cases}
\begin{enumerate}
\item \textbf{Strongly convex $f$ (parameter $\mu>0$).}
   Adding the inequality
   \[
   f(w_{k})-f(w^{\star})\ge \frac{\mu}{2}\|w_{k}-w^{\star}\|^{2},
   \]
   to the derivation yields a linear convergence factor
   \(\bigl(1-\frac{\mu\alpha}{G\sqrt{T}}\bigr)^{T}\) after a suitable
   choice of \(\alpha\).  The detailed algebra follows the same steps as
   above with the extra strong‑convex term.

\item \textbf{Coordinate‑wise AdaGrad (vector case).}
   The scalar analysis applied component‑wise gives
   \[
   f(\bar w_{T})-f(w^{\star})
   \le
   \sum_{j=1}^{d}\frac{D_{j}G_{j}}{\sqrt{T}},
   \]
   where \(D_{j}\) and \(G_{j}\) are the per‑coordinate diameters and
   gradient bounds.

\item \textbf{No bounded‑gradient assumption.}
   If only \(\|g_{k}\|^{2}\le G_{k}\) (i.e. the cumulative sum itself) the bound
   becomes
   \[
   f(\bar w_{T})-f(w^{\star})\le
   \frac{D\sqrt{G_{T}}}{\alpha T}+\frac{\alpha G_{T}}{2T},
   \]
   which still decays as \(O(1/\sqrt{T})\) because \(G_{T}=O(T)\) for
   Lipschitz gradients.
\end{enumerate}

\end{document}